% =========================================================
%  sections/10-referensi.tex
%  Bagian 10 — Daftar Referensi
% =========================================================

\section{Daftar Referensi}
\label{sec:referensi}

Seluruh rujukan yang digunakan dalam penyusunan materi dan asesmen mata kuliah
\NamaMK\ tercantum di bawah ini. Nomor rujukan [R1]--[R14] adalah referensi
utama yang digunakan mahasiswa, sedangkan [R15]--[R35] adalah dokumen RPS,
materi kuliah, dan studi kasus dari perguruan tinggi Indonesia yang menjadi
bahan pembanding dalam penyusunan RPS ini. Seluruh tautan ditelusuri pada
September 2026 dan perlu diverifikasi ulang ketika RPS diperbarui karena isi
laman dapat berubah.

\subsection{Referensi Utama}

\begin{sloppypar}
\begin{enumerate}[label={[R\arabic*]}, leftmargin=3.2em, itemsep=4pt, topsep=4pt]
  \item A. Sweigart, \istilah{Making Games with Python \& Pygame}, 2012.
        \url{https://inventwithpython.com/makinggames.pdf}

  \item A. Sweigart, \istilah{Invent Your Own Computer Games with Python},
        edisi ke-4, 2018.
        \url{https://inventwithpython.com/inventwithpython.pdf}

  \item Pygame Community, \istilah{Pygame Documentation --- Quick Start dan
        Referensi API}. \url{https://www.pygame.org/docs/}

  \item Pygame Community, \istilah{Pygame Documentation --- Sprite} dan
        \istilah{Mask}. \url{https://www.pygame.org/docs/ref/sprite.html} dan
        \url{https://www.pygame.org/docs/ref/mask.html}

  \item R. Nystrom, \istilah{Game Programming Patterns} --- bab
        \istilah{Game Loop}.
        \url{https://gameprogrammingpatterns.com/game-loop.html}

  \item G. Fiedler, ``Fix Your Timestep!'', \istilah{Gaffer On Games}, 2004.
        \url{https://gafferongames.com/post/fix_your_timestep/}

  \item E. Catto, \istilah{Numerical Integration}, GDC 2009.
        \url{https://box2d.org/files/ErinCatto_NumericalIntegration_GDC2009.pdf}

  \item E. dos Santos, \istilah{Game AI: Pathfinding}, bahan kuliah, 2020.
        \url{https://edirlei.com/aulas/game-ai-2020/GAME_AI_Lecture_02_Pathfinding_2020.pdf}

  \item E. dos Santos, \istilah{Game AI: Finite State Machines}, bahan kuliah,
        2020.
        \url{https://edirlei.com/aulas/game-ai-2020/GAME_AI_Lecture_03_Finite_State_Machines_2020.pdf}

  \item E. Bendersky, ``Writing a Game in Python with Pygame, Part III:
        Pathfinding'', 2009.
        \url{https://eli.thegreenplace.net/2009/01/09/writing-a-game-in-python-with-pygame-part-iii}

  \item R. Munir, \istilah{Penentuan Rute, Bagian 2: Algoritma A*}, bahan
        kuliah Institut Teknologi Bandung.
        \url{https://informatika.stei.itb.ac.id/~rinaldi.munir/Stmik/2020-2021/Route-Planning-Bagian2-2021.pdf}

  \item Unity Technologies, \istilah{Introduction to Level Design in Game
        Development and in Unity}, e-book resmi.
        \url{https://assets2.brandfolder.io/bf-boulder-prod/khmctn58vjqkr42wmm3mqvq9/v/1131033155/original/UNITY_Ebook_Introduction%20to%20level%20design%20in%20game%20development%20and%20in%20Unity_V2.pdf}

  \item Unity Technologies, \istilah{Unity Learn --- Unity Essentials} dan
        \istilah{Junior Programmer: Create with Code}.
        \url{https://learn.unity.com/}

  \item Unity Technologies, \istilah{Unity Manual --- Importing Assets} dan
        \istilah{Audio}.
        \url{https://docs.unity3d.com/Manual/ImportingAssets.html}
\end{enumerate}
\end{sloppypar}

\subsection{Referensi Pendukung: RPS dan Materi Perguruan Tinggi Indonesia}

\begin{sloppypar}
\begin{enumerate}[label={[R\arabic*]}, leftmargin=3.2em, itemsep=4pt, topsep=4pt, start=15]
  \item Institut Teknologi Sepuluh Nopember, \istilah{RPS Desain dan
        Pemrograman Game}, Program Studi Teknik Komputer.
        \url{https://www.its.ac.id/komputer/wp-content/uploads/sites/28/2018/03/Desain-dan-pemograman-Game.pdf}

  \item Universitas Negeri Manado, \istilah{RPS Pengembangan Game}, Program
        Studi Teknik Informatika, 2025/2026.
        \url{https://ti.unima.ac.id/rps/view/11}

  \item Universitas Negeri Surabaya, \istilah{RPS Pemrograman Game}, S1
        Pendidikan Teknologi Informasi.
        \url{https://sidia.unesa.ac.id/rps-pdf/s1-pendidikan-teknologi-informasi/pemrograman-game.pdf}

  \item Politeknik Negeri Batam, \istilah{Kurikulum Program Studi Teknologi
        Permainan}, 2024.
        \url{https://www.polibatam.ac.id/wp-content/uploads/2024/06/Kurikulum-Prodi-TP.pdf}

  \item Universitas Abdurrab, \istilah{Deskripsi Mata Kuliah Pemrograman Game},
        SPADA. \url{http://spada.univrab.ac.id/course/info.php?id=1303&lang=en}

  \item Universitas Medan Area, \istilah{Panduan Membuat Game 2D Sederhana
        dengan Pygame}, 2025.
        \url{https://p2dpm.uma.ac.id/2025/10/15/panduan-membuat-game-2d-sederhana-dengan-pygame/}

  \item Universitas Brawijaya, ``Game Development with Python: Let's Build a
        Simple Game!'', \istilah{Gramaswara}.
        \url{https://gramaswara.ub.ac.id/index.php/gramaswara/article/download/185/128}

  \item Universitas Dian Nuswantoro, ``Pemrograman dan Game dengan
        Python/Pygame'', \istilah{Abdimasku}.
        \url{https://abdimasku.lppm.dinus.ac.id/index.php/jurnalabdimasku/article/download/2965/735}

  \item Institut Teknologi Bandung, ``Collision Detection in Precision
        Parry-Based Combat'', makalah IF2123, 2024.
        \url{https://informatika.stei.itb.ac.id/~rinaldi.munir/AljabarGeometri/2024-2025/Makalah/Makalah-IF2123-Algeo-2024%20(31).pdf}

  \item STMIK Jayakarta, ``Penerapan Library Pygame dalam Game RPG The
        Adventure'', \istilah{Jurnal JMI Jayakarta}.
        \url{https://journal.stmikjayakarta.ac.id/index.php/JMIJayakarta/article/download/1013/677/}

  \item Institut Teknologi Nasional Malang, ``Penerapan A* Pathfinding dan
        Finite State Machine pada Game Lost Civilization'', \istilah{JATI}.
        \url{https://ejournal.itn.ac.id/index.php/jati/article/download/5402/3593}

  \item Universitas Muhammadiyah Riau, ``Penerapan Algoritma A* dan Behaviour
        Trees untuk NPC'', \istilah{COSCITECH}.
        \url{https://ejurnal.umri.ac.id/index.php/coscitech/article/download/5418/2453/17612}

  \item UIN Maulana Malik Ibrahim Malang, \istilah{Implementasi FSM dan Logika
        Fuzzy untuk Perilaku NPC}, skripsi.
        \url{http://etheses.uin-malang.ac.id/8204/1/11650053.pdf}

  \item Universitas Negeri Yogyakarta, \istilah{Pengembangan Permainan Edukasi
        KATELU dengan Unity 3D}, skripsi.
        \url{https://eprints.uny.ac.id/52540/1/TAS_AFIF%20FATCHUR%20RAHMAN_12520241062.pdf}

  \item Universitas Diponegoro, ``Media Pembelajaran Budaya Indonesia
        Menggunakan Unity'', \istilah{Jurnal Teknologi dan Sistem Komputer}.
        \url{https://jtsiskom.undip.ac.id/article/download/12667/12286}

  \item Universitas Tanjungpura, ``Game Edukasi Budaya Indonesia'',
        \istilah{JCSKom}.
        \url{https://jurnal.untan.ac.id/index.php/jcskommipa/article/download/49393/75676592053}

  \item STKIP Rokania, \istilah{Pelatihan Dasar Pemrograman Game dengan Unity},
        \istilah{Jurnal MNR}.
        \url{https://e-jurnal.rokania.ac.id/index.php/jmnr/article/download/39/28/172}

  \item Pygame Community, \istilah{Pygame Tutorial: Line-by-Line Chimp}.
        \url{https://www.pygame.org/docs/tut/ChimpLineByLine.html}

  \item Reintech, ``How to Create a Video Game AI with Python''.
        \url{https://reintech.io/blog/how-to-create-a-video-game-ai-with-python}

  \item Institut Teknologi Bandung, ``Implementasi Algoritma A* dalam Navigasi
        Karakter'', makalah IF2120, 2023.
        \url{https://informatika.stei.itb.ac.id/~rinaldi.munir/Matdis/2022-2023/Makalah2022/Makalah-Matdis-2022%20(86).pdf}

  \item STMIK Dharmapala Riau, ``Rancang Bangun Pathfinding A* pada NPC
        Menggunakan Unity Engine'', \istilah{Jurnal JIKB}.
        \url{https://ojs.stmikdharmapalariau.ac.id/index.php/jikb/article/download/831/574}
\end{enumerate}
\end{sloppypar}

\subsection{Catatan Penggunaan Referensi}

\begin{itemize}
  \item Referensi [R1]--[R4] dan [R32] digunakan sebagai acuan utama praktikum
        CPMK1 dan CPMK2 karena membahas langsung API Pygame.
  \item Referensi [R5]--[R7] digunakan sebagai landasan teori \istilah{game
        loop}, \istilah{timestep}, dan integrasi numerik.
  \item Referensi [R8]--[R11] dan [R34] digunakan sebagai landasan teori FSM,
        A*, serta \istilah{pathfinding} pada grid.
  \item Referensi [R12]--[R14] digunakan sebagai acuan teknis resmi untuk
        CPMK3 dan CPMK4.
  \item Referensi [R15]--[R35] digunakan sebagai pembanding struktur RPS,
        cakupan materi, dan studi kasus penerapan pada perguruan tinggi
        Indonesia; dokumen penelitian dipakai sebagai studi kasus, bukan
        pengganti dokumentasi teknis resmi.
\end{itemize}
